%% ======================================================================
%%  Exercises (pp. 97-105)
%% ======================================================================
\clearpage
\folio{97}%
\global\let\currentseccode\empty
\global\let\pagefirstsec\empty
\setcounter{equation}{0}%
\renewcommand{\chaptertitle}{}%
\def\currentsectitle{}%
\thispagestyle{emptyhead}%
\phantomsection
\pdfbookmark[0]{Exercises}{fm:exercises}%
\label{fm:exercises}%
\addtocontents{toc}{\string\fml{EXERCISES}{\thepage}{fm:exercises}}%
\begin{center}{\Large Exercises}\end{center}%
\vspace{1.2\baselineskip}%
\par

The exercises marked A are easy exercises which the reader may use to check
his comprehension of the text. Those marked B are more elaborate. The review
problems at the end have been given as examination problems in Paris.

\newcounter{exno}
\newcommand{\exchap}[1]{\def\exprefix{#1}\setcounter{exno}{0}}
\newcommand{\exitem}[1]{\par\medskip\refstepcounter{exno}%
  \noindent{\bfseries #1.}\label{exer:\exprefix.\arabic{exno}}\quad}
\exchap{I}

\divhead{Chapter I}

\exitem{1B} Let $p$ be a prime number and let $r\in\mathbf{N}$. Show that
the multiplicative group $(\mathbf{Z}/p^{r}\mathbf{Z})^{*}$ is cyclic except
if $p=2$ and $r\geq3$ (for $p$ odd, show that the residue class of $1+p$ is
of order $p^{r-1}$; for $p=2$ look at the order of the residue class of $5$;
then use Corollary \ref{cor:I.5.4} of Theorem \ref{thm:I.5.1}, \hyperref[sec:1.5]{§5}).

\exitem{2B} \begin{alist}
\item Show that the diophantine equation $X^{2}+Y^{2}=3Z^{2}$ has
no non-trivial solution (i.e.\ $\neq(0,0,0)$). (Reduce mod $3$.) Show the
same for $X^{2}+Y^{2}=7Z^{2}$ and $X^{2}+Y^{2}=11Z^{2}$.
\item Show that $X^{2}+Y^{2}=5Z^{2}$ and $X^{2}+Y^{2}=13Z^{2}$ have
non-trivial integer solutions.
\item Try to generalize (a) to certain equations $X^{2}+Y^{2}=pZ^{2}$, $p$ a
given prime (study under what condition $-1$ is a square in
$\mathbf{F}_{p}^{*}$).
\end{alist}

\exitem{3B} Rediscover the classical proof for the fact that there are
infinitely many prime numbers. By analogous methods show that there are
infinitely many primes of the form $4k-1$ ($k\in\mathbf{N}$).

\exitem{4B} In order that $n\in\mathbf{N}$ be prime it is necessary and
sufficient that $n$ divide $(n-1)!+1$. (If $n$ is prime, calculate the
product of the elements of $\mathbf{F}_{n}^{*}$; check what happens in the
case when $n$ is not prime.)

\exitem{5B} Show that in a finite field $K$ every element is a sum of two
squares (first treat the case where $q=\operatorname{card}(K)$ is even; if
$q$ is odd, calculate the number of values taken by the functions
$x\mapsto x^{2}$ and $y\mapsto a-y^{2}$, $x,y\in K$, $a\in K$ given).

\exitem{6B} Factor the polynomial $X^{3}-X+1$ in $\mathbf{F}_{23}$ and the
polynomial $X^{3}+X+1$ in $\mathbf{F}_{31}$ (each has a double root and a
simple root).

\exitem{7A} Give an example of two ideals $\mathfrak{a},\mathfrak{b}$ of a
ring $A$ such that
$\mathfrak{a}\cap\mathfrak{b}\neq\mathfrak{a}\mathfrak{b}$. Show that it is
always the case that
$\mathfrak{a}\mathfrak{b}\subset\mathfrak{a}\cap\mathfrak{b}$.

\exitem{8B} Let $A$ be an integral domain, $a$ and $b\in A$, and let
$B=A[X]/(aX+b)$. Show that, if
$Aa\cap Ab=Aab$, then $B$ is an integral domain (consider the element
$-b/a$ of the field of fractions $K$ of $A$, and show that the homomorphism
$\varphi:A[X]\to K$ defined by $\varphi(X)=-b/a$ and $\varphi(y)=y$ for
$y\in A$ has precisely $(aX+b)$ as its kernel).

\exitem{9B} Let $A$ be an integral domain and let $i$ and $j$ be relatively
prime positive integers. Show that the ideal $(X^{i}-Y^{j})$ in the
polynomial ring $A[X,Y]$ is prime (define a homomorphism
$A[X,Y]\to A[X]$ having that ideal as its kernel).

\exitem{10B} (a) Let $A$ be an integral domain, $K$ its field of fractions,
and $b$ a non-zero element of $A$. Show that the $A$-module
$\operatorname{Hom}(A/Ab,K/A)$ is isomorphic to $A/Ab$ (recall that, given
two $A$-modules $E$ and $F$, we write $\operatorname{Hom}(E,F)$ for the set
of homomorphisms, or $A$-linear maps, from $E$ to $F$; recall that
$\operatorname{Hom}(E,F)$ is an $A$-module. We note that a homomorphism
$\varphi:A/Ab\to K/A$ is determined by its value at the residue class of
$1$).
\begin{alist}
\item For an $A$-module $E$, put $E'=\operatorname{Hom}(E,K/A)$ and put
$E''=\operatorname{Hom}(E',K/A)$. Define a canonical homomorphism
$E\to E''$.
\item Suppose that $A$ is a principal ideal ring. Let $E$ be a torsion
module of finite type over $A$ (i.e.\ such that, for every $x\in E$, there
exists non-zero $a\in A$ for which $ax=0$). Determine the structure of $E$
and deduce from (a) that $E$ is isomorphic to
$E'=\operatorname{Hom}(E,K/A)$. Deduce from this that the canonical
homomorphism $E\to E''$ is an isomorphism.
\end{alist}

\exchap{II}

\divhead{Chapter II}

\exitem{1A} Theorem \ref{thm:II.7.1} of \hyperref[sec:2.7]{§7} remains true if, instead of
supposing that $K$ is of characteristic zero, we suppose $K$ finite. Tell
why. Is the case where $K$ is finite of any interest?

\exitem{2B} Let $A$ be an integrally closed ring, $K$ its field of
fractions, and $P(X)\in A[X]$ a monic polynomial. If $P(X)$ is reducible in
$K[X]$, show that it is reducible in $A[X]$ (consider the roots of $P(X)$ in
an extension of $K$).

\exitem{3B} For every positive integer $n$ write $f(n)$ for the number of
monic irreducible polynomials of degree $n$ over $\mathbf{F}_{q}$. Show that
$\sum_{d\mid n}df(d)=q^{n}$ (classify the elements of $\mathbf{F}_{q}^{n}$
according to their minimal polynomials over $\mathbf{F}_{q}$, and note that
$\mathbf{F}_{q}^{d}\subset\mathbf{F}_{q}^{n}$ if and only if $d\mid n$).
Deduce from this the values of $f(1),f(2),f(3),f(4)$, and $f(p)$ for $p$
prime.

\exitem{4A} Give an example of an integral domain which is not integrally
closed (look at \hyperref[sec:2.5]{§5}).

\exitem{5B} Let $R$ be a ring, $A$ a subring of $R$, and $x$ a unit in $R$.
Show that every $y\in A[x]\cap A[x^{-1}]$ is integral over $A$ (show that
there exists an integer $n$ such that the $A$-module
$M=A+Ax+\cdots+Ax^{n}$ is stable under multiplication by $y$, i.e.\ that
$yM\subset M$. One will at this point find the idea needed to conclude the
proof by looking at the last part of the proof of Theorem
\ref{thm:II.1.1}, \hyperref[sec:2.1]{§1}).

\exchap{III}

\divhead{Chapter III}

\exitem{1A} Why is a field a Dedekind ring?

\exitem{2A} How could one modify the statement of Theorem \ref{thm:III.4.2},
\hyperref[sec:3.4]{§4} so that the hypothesis that $A$ is not a field would become superfluous?

\exitem{3B} In a Noetherian ring $A$ it is clear that any ideal different
from $A$ itself is contained in a maximal ideal. Show that this is still
true when $A$ is not Noetherian (use Zorn's lemma).

\exitem{4A} Give an example of an integral domain $A$ containing a prime
ideal $\mathfrak{p}\neq0$ and a subring $A'$ such that
$\mathfrak{p}\cap A'=(0)$.

\exitem{5A} Let $K$ be a field. Is the ring $K[X,Y]$ of polynomials in two
variables over $K$ a Dedekind ring?

\exitem{6} Let $K$ be a field and let $A$ be the subring
$K[X^{2},X^{3}]$ of the ring of polynomials $K[X]$.
\begin{alist}
\item By using the inclusion relations
$K[X^{2}]\subset A\subset K[X]$, show that $A$ is Noetherian and that every
non-zero prime ideal is maximal.
\item Show that $A$ is not a Dedekind ring (consider the element $X$).
\end{alist}

\exchap{IV}

\divhead{Chapter IV}

\exitem{1B} Show that the ring of integers $A$ of the cubic field
$\mathbf{Q}[x]$ with $x^{3}=2$ is principal (majorize the discriminant of
$A$ by $D(1,x,x^{2})$ and apply Corollary \ref{cor:IV.3.1} of Proposition
\ref{prop:IV.3.1}, \hyperref[sec:4.3]{§3}; in fact, $A=\mathbf{Z}[x]$. See a review problem below.)

\exitem{2B} Determine the units of the ring $A=\mathbf{Z}[X]/(X^{3})$ and the
structure of the group $A^{*}$.

\exitem{3B} Let $K$ be a cubic field such that $r_{1}=r_{2}=1$. Suppose $K$
is imbedded in $\mathbf{R}$.
\begin{alist}
\item Show that the positive units of $K$ form a group isomorphic to
$\mathbf{Z}$, and that every positive unit of $K$ is of norm $1$.
\item Let $d$ be the absolute discriminant of $K$ and let $u$ be a unit
greater than $1$. Show that $|d|\leq4u^{3}+24$ (put $u=x^{2}$ with
$x\in\mathbf{R}$ and $x>0$. Note that the conjugates of $u$ are of the form
$x^{-1}e^{iy}$ and $x^{-1}e^{-iy}$ with $y\in\mathbf{R}$. Calculate the
discriminant $d'=D(1,u,u^{2})$ as a function of $x$ and $y$, say
$|d'|^{1/2}=\varphi(x,y)$. Find the minimum of $\varphi(x,y)$ for fixed $x$
and deduce that $|d'|\leq4u^{3}+24$. Conclude by observing that $d$ divides
$d'$).
\item Show that the polynomial $X^{3}+10X+1$ is irreducible over
$\mathbf{Q}$ (cf.\ \hyperref[sec:5.3]{Chapter V, §3}, example near the end). Let $x$ be one of
its real roots. Show that the ring of integers of the cubic field
$\mathbf{Q}[x]$ is $\mathbf{Z}[x]$ (ibid.). Show that $u=-1/x$ is a
generator of the group of positive units of $\mathbf{Q}[x]$.
\end{alist}

\exitem{4B} Let $K$ be a number field, let $A$ be the ring of integers of
$K$, and let $P$ denote the set of maximal ideals of $A$. For
$\mathfrak{p}\in P$ and $x\in K^{*}$ write
$v_{\mathfrak{p}}(x)$ for the exponent of $\mathfrak{p}$ in the
factorization of $Ax$ into a product of prime ideals. Put
$v_{\mathfrak{p}}(0)=+\infty$.
\begin{alist}
\item Show that, for $x,y\in K$,
$v_{\mathfrak{p}}(xy)=v_{\mathfrak{p}}(x)+v_{\mathfrak{p}}(y)$ and
$v_{\mathfrak{p}}(x+y)\geq\min(v_{\mathfrak{p}}(x),v_{\mathfrak{p}}(y))$.
\item Let $P'\subset P$. Show that the set $A(P')$ of elements $x\in K$
such that $v_{\mathfrak{p}}(x)\geq0$ for all $\mathfrak{p}\in P'$ is a
subring of $K$. Consider the cases $P'=P$ and $P'=\phi$.
\item Take for $P'$ the complement of a finite set $S\subset P$. Show that
the group $U$ of units of $A(P')$ is of finite type and that $U/A^{*}$ is a
free $\mathbf{Z}$-module of rank $\operatorname{card}(S)$ (consider the
mapping $x\mapsto(v_{\mathfrak{p}}(x))_{\mathfrak{p}\in S}$ of $U$ to
$\mathbf{Z}^{S}$; find its kernel; determine its image by observing that
every ideal of $A$ has a power which is a principal ideal).
\end{alist}

\exchap{V}

\divhead{Chapter V}

\exitem{1B} Let $A$ be a principal ideal ring and let $K$ be its field of
fractions.
\begin{alist}
\item Show that every intermediate ring $B$ ($A\subset B\subset K$) is a
ring of fractions of $A$ (if $a/b\in B$ with $a,b\in A$ relatively prime,
then show, by using Bezout's identity, that $1/b\in B$).
\item Classify the rings of fractions of $A$ (consider the prime elements of
$A$ which are invertible in $S^{-1}A$).
\item Show that any ring of fractions of $A$ is principal.
\end{alist}

\exitem{2A} Show that no integer $\equiv7$ modulo $8$ is a sum of three
squares (reduce mod $8$).

\exitem{3B} Let $K$ be a field and $A$ a reduced $K$-algebra of finite
dimension over $K$.
\begin{alist}
\item Making use of \hyperref[sec:5.3]{§3}, Lemma \ref{lem:I.3.2}, show that $A$ is isomorphic
to a finite product $\prod_{i=1}^{n}L_{i}$ of extensions of finite degree of
$K$.
\item Show that, if $K$ is of characteristic $0$, $A$ has a primitive
generator, i.e.\ $A$ is of the form $K[x]$ with $x\in A$ (find a set of
primitive generators $(x_{i})$ for the fields $(L_{i})$ over $K$ whose
minimal polynomials over $K$ are pairwise distinct. Take
$x=(x_{1},\ldots,x_{n})$).
\end{alist}

\exitem{4B} Let $p$ be a prime number $\equiv1\bmod4$. Show that there are
exactly $8$ pairs $(a,b)\in\mathbf{Z}\times\mathbf{Z}$ such that
$a^{2}+b^{2}=p$ (in other words, the diophantine equation $p=a^{2}+b^{2}$
has only one solution up to trivial transformations).

\exitem{5B} Let $K$ be a field, $A$ an algebra of finite dimension over $K$.
Give examples where $A$ has no primitive generator (i.e.\ not of the form
$K[x]$):
\begin{alist}
\item When $A$ is not reduced;
\item When $A$ is reduced and $K$ is finite (cf.\ \hyperref[exer:V.3]{Exercise 3}).
\end{alist}

\exitem{6B} (a) Show that the polynomial $X^{5}-X+1$ is irreducible over
$\mathbf{Q}$ (look at its reduction mod $5$). Let $x$ be one of its roots.
Calculate the integers $r_{1}$ and $r_{2}$ for the field $\mathbf{Q}[x]$.
\begin{alist}
\item Calculate the discriminant of $(1,x,\ldots,x^{4})$ (over
$\mathbf{Z}$; cf.\ \hyperref[sec:2.7]{Chapter II, §7}, formula \eqref{eq:II.7.7}). Show that it
is square-free and deduce that $\mathbf{Z}[x]$ is the ring of integers of
$\mathbf{Q}[x]$ (\hyperref[sec:5.3]{§3}, Proposition \ref{prop:IV.3.1}).
\item Show that the ring $\mathbf{Z}[x]$ is principal (use Corollary
\ref{cor:IV.3.1} of Proposition \ref{prop:IV.3.1} of \hyperref[sec:4.3]{§3, Chapter IV} and note,
by reduction of $X^{5}-X+1$ mod $2$ and mod $3$, that $\mathbf{Z}[x]$
contains no ideals of norm $2$ or $3$).
\end{alist}

\exitem{7B} Let $K=\mathbf{Q}[\sqrt{d}]$ be a quadratic field where $d$ is a
square-free integer and let $A$ be the ring of integers of $K$.
\begin{alist}
\item By means of Corollary \ref{cor:IV.3.1}, \hyperref[sec:4.3]{§3, Chapter IV}, show that $A$
is principal for $d=2,3,5,13$,
$-1,-2,-3,-7$.
\item By the same method show that any ideal of $K$ is equivalent to an
integral ideal of norm $1$ or $2$ for $d=6,7,17,21,29,33$,
$-5,-11,-15,-19$.
In which of the preceding fields does $2$ remain prime
($d=21,29,-11,-19$; cf.\ \hyperref[sec:5.4]{§4}, \hyperref[prop:V.4.1]{Proposition 1})? Show that their rings of
integers are principal.
\item Suppose $d\equiv2$ or $3\bmod4$. Show that $K$ contains a principal
ideal of norm $2$ if and only if there exist $a$ and $b\in\mathbf{Z}$ such
that $a^{2}-db^{2}=\pm2$. Deduce that $A$ is principal for $d=6,7$, and that
it contains two ideal classes when $d=-5$ (note that $2$ ramifies in this
field).
\item Suppose $d\equiv1\bmod4$. Show that $K$ contains a principal ideal of
norm $2$ if and only if $a^{2}-db^{2}=\pm8$ has a solution
$(a,b)\in\mathbf{Z}\times\mathbf{Z}$. Show that $A$ is principal for
$d=17,33$, and that it contains two ideal classes when $d=-15$ (note that
$2$ splits in this field).
\item Show that, for $d=10$ and $d=-6$, $A$ contains two classes of ideals
(note that any ideal of $K$ is equivalent to an integral ideal of norm
$1,2$, or $3$, study the factorizations of $2$ and $3$ in $K$, and show that
$(2+\sqrt{10})(2-\sqrt{10})=-2\cdot3$, also
$(\sqrt{-6})^{2}=-2\cdot3$).
\item Show that, for $d=-23$, $A$ contains three classes of ideals (proceed
as in (c)). Show that there are non-principal prime ideals
$\mathfrak{p},\mathfrak{p}',\mathfrak{q},\mathfrak{q}'$ of $A$ with the
property that $2A=\mathfrak{p}\mathfrak{p}'$ and
$3A=\mathfrak{q}\mathfrak{q}'$. Observe that
$x=\frac{1}{2}(3+\sqrt{-23})$ and $y=\frac{1}{2}(1+\sqrt{-23})$ are elements
of $A$ with norms $8$ and $6$, respectively. Study the factorization of $Ax$
and $Ay$ into prime ideals).
\end{alist}

\exitem{8A} Let $K,K'$ and $K''$ be number fields with
$K\subset K'\subset K''$. Let $\mathfrak{p}''$ be a maximal ideal of $K''$
and let $\mathfrak{p}'=\mathfrak{p}''\cap K'$. Write $e$ (respectively,
$e'$, $e''$) for the ramification index of $\mathfrak{p}''$ over $K$
(respectively, $\mathfrak{p}'$ over $K$, $\mathfrak{p}''$ over $K'$). Write
$f$ (respectively, $f'$ and $f''$) for the residual degree of
$\mathfrak{p}''$ over $K$ (respectively, $\mathfrak{p}'$ over $K$,
$\mathfrak{p}''$ over $K'$). Show that $e=e'e''$ and $f=f'f''$.

\exitem{9B} (a) Show that the polynomial $X^{3}+X^{2}-2X+8$ is irreducible
over $\mathbf{Q}$ (note that it has no root in $\mathbf{Z}$, so no root in
$\mathbf{Q}$ either). Write $x$ for one of its roots, $K$ for the cubic
field $\mathbf{Q}[x]$, and $A$ for the ring of integers of $K$.
\begin{alist}
\item Show that $D(1,x,x^{2})=4503$ and $503$ is prime.
\item Show that $y=4/x$ is an integer of $K$, that $y\notin\mathbf{Z}[x]$,
and that $A'=\mathbf{Z}+\mathbf{Z}x+\mathbf{Z}y$ is a subring of $K$ which
contains $\mathbf{Z}[x]$ properly. Show that the discriminant of $A'$ over
$\mathbf{Z}$ is $503$ (compare the discriminants of $A$ and $A'$, or
calculate $D(1,x,y)$ directly). Conclude that $A=A'$.
\item Show that $2A$ is the product of $3$ distinct prime ideals (write the
multiplication table of the elements $1,x,y$ and reduce mod $2$). Deduce
from this fact that there exists no element $u$ of $A$ such that
$A=\mathbf{Z}[u]$ (note that $A/2A$ is not primitively generated over
$\mathbf{F}_{2}$).
\item Show that, if $u$ is a primitive element of $K$ contained in $A$, the
group $A/\mathbf{Z}[u]$ admits a quotient group of order $2$ (note that two
of the three homomorphisms of $A$ on $\mathbf{F}_{2}$ given by (d) coincide
on $\mathbf{Z}[u]$, and consider their difference). Conclude that
$D(1,u,u^{2})$ is a multiple of $4\cdot503$ (note that
$D(1,u,u^{2})=503\operatorname{card}((A/\mathbf{Z}[u])^{2})$). Thus none of
the numbers $D(1,u,u^{2})$ generates the discriminant of $K$ as an ideal in
$\mathbf{Q}$.
\item Show that the ring $A$ is principal (by \hyperref[sec:4.3]{Chapter IV, §3}, \hyperref[cor:IV.3.1]{Corollary 1} to
\hyperref[prop:IV.3.1]{Proposition 1} it suffices to show that every prime ideal of norm $n\leq6$
is principal; exclude the case $n=6$ and, by making use of (d), the case
$n=4$; then study the norms and the decomposition of $x,y,x-1,x+2$, and
$x+3$ in $A$; in order to settle the case $n=2$, affirm that $(x-2)/(x+2)$
and $(x-1)/(x+3)$ generate two of the three prime ideals of norm $2$; for
$n=3$, show that $3A$ is prime; to settle $n=5$, show that $5A$ is the
product of a prime ideal of norm $5$ and of another of norm $25$, and
compare the factorizations and norms of $5A$ and of $(x+1)A$).
\end{alist}

\exitem{10B} (a) Let $p$ and $q$ be distinct prime numbers $\geq5$ and such
that $pq^{2}\not\equiv1\pmod9$. Let $K$ be the cubic field $\mathbf{Q}[u]$
where $u^{3}=p^{2}q$. Show that $u$ and $v=pqu^{-1}$ are integers of $K$,
that $D(1,u,v)=-27p^{2}q^{2}$, and that
$B=\mathbf{Z}+\mathbf{Z}u+\mathbf{Z}v$ is a subring of the ring $A$ of
integers of $K$ (note that $u^{2}=qv$ and $v^{2}=pu$).
\begin{alist}
\item Let $\mathfrak{q}$ be a prime ideal of $A$ containing $Aq$. Deduce
from the fact that $u^{3}=p^{2}q$ that the exponent of $\mathfrak{q}$ in the
decomposition of $Aq$ is $3$, and that $Aq=\mathfrak{q}^{3}$. Show that
$B\cap\mathfrak{q}$ is the ideal of $B$ generated by
$u,v,q$, that $A/\mathfrak{q}=B/B\cap\mathfrak{q}$, and that
$A=B+\mathfrak{q}$. Conclude from this that $A=B+\mathfrak{q}A$. Similarly
show that $Ap=\mathfrak{p}^{3}$ with $\mathfrak{p}$ prime and that
$A=B+\mathfrak{p}A$.
\item By considering $(u\pm1)^{3}$ (which differs by one from the class of
$p^{2}q\bmod3$), show that $3A$ is the cube of a prime ideal of $A$, and
that (as in (b)) $A=B+3A$.
\item Let $x=a+bu+cv$ ($a,b,c\in\mathbf{Q}$), an integer of $K$. Deduce, by
calculating traces, that $3pqx\in B$, from which it follows that
$3pqA\subset B$. Use (b) and (c) to show that $A=B+3pqA$. Conclude from
this that $A=B$.
\item Show that, for $x=a+bu+cv$ ($a,b,c\in\mathbf{Z}$),
\[
D(1,x,x^{2})=3p^{2}q^{2}\left(b^{3}q-c^{3}p\right)^{2}.
\]
\item Deduce from (e) that, if $p\equiv1\pmod3$ and
$q^{(p-1)/3}\equiv1\pmod p$, $A$ has no base (over $\mathbf{Z}$) of the form
$(1,x,x^{2})$ ($x\in A$). Show that this is the case for $p=7$ and $q=5$
(in which case $K=\mathbf{Q}[\sqrt[3]{175}]=\mathbf{Q}[\sqrt[3]{245}]$).
Show that, in this case, the absolute discriminant of $K$ is the greatest
common divisor of the numbers $D(1,x,x^{2})$, for $x\in A$.
\end{alist}

\exchap{VI}

\divhead{Chapter VI}

\exitem{1B} (a) Show that the field
$L=\mathbf{Q}[\sqrt{5},\sqrt{-1}]$ is a Galois extension of degree $4$ of
$\mathbf{Q}$. Determine its Galois group.
\begin{alist}
\item Show that the ring $A$ of integers of $L$ is
$\mathbf{Z}[\sqrt{-1},(1+\sqrt{5})/2]$ (note that the discriminant of
$(1,(1+\sqrt{5})/2)$ over the ring of Gaussian integers is square-free
(i.e.\ no Gaussian integer square factors); then make use of ideas from
\hyperref[sec:5.3]{Chapter V, §3}, \hyperref[prop:V.3.1]{Proposition 1}). Calculate the absolute discriminant of $L$.
Show that the only primes which ramify in $L$ are $2$ and $5$, and that the
corresponding ramification indices are both $2$ (use the fact that $n=efg$,
cf.\ \hyperref[sec:6.2]{§2}, \hyperref[prop:VI.2.1]{Proposition 1}).
\item Calculate the Frobenius automorphisms
$\left(\frac{L/\mathbf{Q}}{p}\right)$ for $p$ prime distinct from $2$ and
$5$ (cf.\ \hyperref[sec:6.3]{§3}). Calculate the decomposition and inertia groups for $2$ and
$5$.
\item Show that $\mathbf{Q}[\sqrt{-5}]\subset L$, and that no prime ideal of
$\mathbf{Q}[\sqrt{-5}]$ ramifies in $L$.
\item Similarly, show that
$\mathbf{Q}[\sqrt{10}]\subset\mathbf{Q}[\sqrt{2},\sqrt{5}]$, and that no
prime ideal of $\mathbf{Q}[\sqrt{10}]$ ramifies in
$\mathbf{Q}[\sqrt{2},\sqrt{5}]$.
\end{alist}

Note that both $K=\mathbf{Q}[\sqrt{-5}]$ and $K=\mathbf{Q}[\sqrt{10}]$, by
the above, admit abelian ``unramified'' extensions (i.e.\ in which no prime
ideal of $K$ ramifies) whose Galois group (cyclic of order $2$) is
isomorphic to the group of ideal classes of $K$. This is a very special case
of the celebrated ``class-field theory''.

\exitem{2B} Let $x$ be the real root of the polynomial $X^{3}-X+1$, $y$ and
$\bar{y}$ its two other roots in $\mathbf{C}$, and $K$ the cubic field
$\mathbf{Q}[x]$ (cf.\ \hyperref[sec:5.3]{Chapter V, §3}, example).
\begin{alist}
\item Show that $y+\bar{y}=-x$, $y\bar{y}=-1/x$, and
$[(y-x)(\bar{y}-x)(y-\bar{y})]^{2}=-23$ (cf.\ \hyperref[sec:2.7]{Chapter II, §7}, formulas
\eqref{eq:II.7.6} and \eqref{eq:II.7.7}). Conclude from this that
$K[\sqrt{-23}]=\mathbf{Q}[x,y,\bar{y}]$ and that this field $L$ is a Galois
extension of degree $6$ of $\mathbf{Q}$. Determine its Galois group $G$
(note that a $\mathbf{Q}$-automorphism of $L$ is determined by its effect on
$x,y,\bar{y}$).
\item What are the subgroups of $G$? Determine the corresponding subfields
of $L$ (cf.\ \hyperref[sec:6.1]{§1}, \hyperref[thm:VI.1.2]{Theorem 2}). What are the subfields of $L$ which are Galois
over $\mathbf{Q}$?
\item Recall that the ring of integers of $K$ is $\mathbf{Z}[x]$ (\hyperref[chap:V]{Chapter V},
the example at the end of \hyperref[sec:5.3]{§3}). Show that it is principal (\hyperref[sec:4.3]{Chapter IV, §3},
\hyperref[cor:IV.3.1]{Corollary 1} of \hyperref[prop:IV.3.1]{Proposition 1}).
\item Show that $23$ is the only prime number ramified in $L$, and that the
corresponding ramification index is $2$ (note that the decomposition of $23$
into prime factors in the ring of integers of $K$ is of the form
$\alpha\beta^{2}$. Denoting the numbers defined in \hyperref[prop:VI.2.1]{Proposition 1} of \hyperref[sec:6.2]{§2}
relative to the decomposition of $23$ in $L$ by $e,f$, and $g$, deduce that
$g\geq2$ and that $e$ is even (cf.\ \hyperref[exer:V.8]{exercise 8} of \hyperref[chap:V]{Chapter V}). Use the fact
that $efg=6$ to conclude that $e=2,f=1,g=3$).
\item Show that $\mathbf{Q}[\sqrt{-23}]\subset L$, and that no prime ideal of
$\mathbf{Q}[\sqrt{-23}]$ ramifies in $L$.
\end{alist}

Observe that $L$ is a cyclic unramified extension of degree $3$ of
$\mathbf{Q}[\sqrt{-23}]$, and that the ideal class group of
$\mathbf{Q}[\sqrt{-23}]$ is cyclic of order $3$ (cf.\ \hyperref[exer:V.7]{Exercise 7}, (f) of
\hyperref[chap:V]{Chapter V}). This is another special case of class field theory.

\exitem{3B} Reformulate \hyperref[exer:VI.2]{exercise 2} to apply to the polynomial $X^{3}+X+1$
(instead of $X^{3}-X+1$) and to $-31$ (in place of $-23$). Reformulate for
$\mathbf{Q}[\sqrt{-31}]$ \hyperref[exer:V.7]{Exercise 7} (f) of \hyperref[chap:V]{Chapter V}.

\exitem{4B} Let $K$ be a number field, and let $L$ be the logarithmic
imbedding of $K^{*}$ in $\mathbf{R}^{+}$ (\hyperref[sec:4.4]{Chapter IV, §4}, \hyperref[thm:IV.4.1]{Theorem 1}). We
write $U$ for the group of units of $K$ and we call the volume of the
lattice $L(U)$ of the hyperplane $W$ considered in \hyperref[sec:4.4]{Chapter IV, §4},
\hyperref[thm:IV.4.1]{Theorem 1} (with respect to the measure on $W$ induced from Lebesgue measure
by the mapping $W\to\mathbf{R}^{r_{1}+r_{2}-1}$) the \emph{regulator} of
$K$.
\begin{alist}
\item Calculate the regulators of the real quadratic fields given as
explicit examples in \hyperref[sec:4.6]{Chapter IV, §6}.
\item Let $z$ be a primitive $n$th root of unity in $\mathbf{C}$, let $K$ be
the cyclotomic field $\mathbf{Q}[z]$, let $\sigma$ be the automorphism of
$K$ defined by $\sigma(z)=z^{-1}$ and $K'$ the field left fixed by
$\sigma$. Calculate $[K:K']$ and the ranks of the groups of units $U$ and
$U'$ of $K$ and $K'$, respectively. Show that these ranks are equal to the
same integer $r$. Show that there exists an integer $m$ such that the
regulators $\rho$ and $\rho'$ of $K$ and $K'$ are connected by the relation
$\rho'=m\rho$. Calculate the number of roots of unity contained in $K$
(respectively, $K'$) (one may take $z=e^{2\pi i/n}$). Determine the index of
$U'$ in $U$.
\item With the notations of (b), suppose that $n=p^{s}$ with $p$ prime,
$p\neq2$. Show that for $u\in U$, we have $\sigma(u)=\pm z^{2t}u$ (note that
$\sigma(u)u^{-1}$ has absolute value one for any imbedding of $K$ in
$\mathbf{C}$, thus is a root of unity), whence $\sigma(z^{t}u)=\pm z^{t}u$.
Show that $\sigma(z^{t}u)=z^{t}u$ (exclude the minus sign by considering
the conjugates of $z^{t}u$ over $K'$). Deduce that $U=VU'$, where $V$ is the
group of $n$th roots of unity contained in $K$. Show that the integer $m$
of (b) is equal to $1$.
\end{alist}

\exitem{5B} (a) Let $K$ be the cyclotomic field generated over
$\mathbf{Q}$ by a ninth root of unity. Show that it is a cyclic extension of
degree $6$ of $\mathbf{Q}$ (cf.\ \hyperref[sec:6.4]{§4}). Let $G$ be its Galois group.
\begin{alist}
\item By making use of the subgroup $H$ of index $3$ of $G$, construct a
Galois extension $K'$ of degree $3$ of $\mathbf{Q}$ (cf.\ \hyperref[sec:6.1]{§1}, \hyperref[thm:VI.1.2]{Theorem 2}).
Show that $K'=\mathbf{Q}[\cos2\pi/9]$ and determine the minimal polynomial
of $2\cos2\pi/9$ over $\mathbf{Q}$ (it has as roots $z+z^{8}$, $z^{2}+z^{7}$,
and $z^{4}+z^{5}$).
\end{alist}

%% ----------------------------------------------------------------------
\divhead{REVIEW PROBLEMS}

\exchap{R1}

\subsechead{}{I}

In the field $\mathbf{C}$ of complex numbers, put $z=e^{2\pi i/10}$ and
write $K$ for the cyclotomic field $\mathbf{Q}[z]$.

\exitem{1} Find the degree, the ring of integers, and the absolute
discriminant of $K$ (to calculate the absolute discriminant it will be
convenient to calculate the numbers $\operatorname{Tr}(z^{i})$).
What primes ramify in $A$?

\exitem{2} Write $\mathbf{R}$ for the field of real numbers and put
$K'=K\cap\mathbf{R}$. Show that $K'$ is a quadratic field and determine
which one. Describe the ring $A'$ of integers of $K'$ and the group
$A'^{*}$ of units of $A'$.

\exitem{3} Show that $K'$ is the only quadratic field contained in $K$
(note that, if $K''$ is a subfield of $K$, every prime number which ramifies
in $K''$ ramifies in $K$).

\exitem{4} Show that $K$ contains no roots of unity besides the tenth roots
of unity (it suffices to consider those roots of unity whose order is prime
or a power of $2$, and to make use of \hyperref[exer:R1.3]{(3)}). Write $G$ for the group of tenth
roots of unity.

\exitem{5} (a) Making use of the logarithmic imbedding of $K$, show that,
if $a$ is a unit of absolute value one in the ring $A$ of integers of $K$,
then $a$ is a root of unity.
\begin{alist}
\item Let $a=re^{i\theta}$ be a unit of $A$ ($r>0$). Show that
$r^{2}\in A'^{*}$ and $e^{2i\theta}\in G$. Also show that the relations
$r\in A'^{*}$ and $e^{i\theta}\in G$ are equivalent.
\item Recall that
$A(1-z)=A(1-z^{2})=A(1-z^{3})=A(1-z^{4})$ is a prime ideal
$\mathfrak{p}$ of $A$ and that $5A=\mathfrak{p}^{4}$. Conclude that no
element of $(1-z^{3})K'$ is a unit of $A$, that $A$ contains no unit of
argument $\pi/10$, and that the group $A^{*}$ of units of $A$ is a direct
product of $G$ and of the group $U$ of positive units of $A'$.
\end{alist}

\exitem{6} Recall that, with the classical notations, every class of ideals
of a number field contains an integral ideal of norm
$\leq(4/\pi)^{r_{2}}n!n^{-n}|d|^{1/2}$. Show that the ring $A$ is principal.

\exchap{R2}

\subsechead{}{II}

Let $\mathbf{Q}$ be the field of rational numbers, $x$ the real number such
that $x^{3}=2$, and $K$ the field $\mathbf{Q}(x)$.

\exitem{1} Show that the polynomial $X^{3}-2$ is irreducible over
$\mathbf{Q}$, and that $K$ is an extension of degree $3$ of $\mathbf{Q}$. In
the field $\mathbf{C}$ of complex numbers what are the fields conjugate to
$K$?

\exitem{2} Let $z=a+bx+cx^{2}$ ($a,b,c\in\mathbf{Q}$) be an element of $K$.
Calculate the trace, norm, and characteristic polynomial of $z$ over
$\mathbf{Q}$.

\exitem{3} Let $\mathbf{Z}$ be the ring of rational integers and $B$ the
subring $\mathbf{Z}[x]$ of $K$. In this problem we intend to show that the
ring $A$ of integers of $K$ is equal to $B$.
\begin{alist}
\item Show that $B\subset A$, and find a base for $B$ as a
$\mathbf{Z}$-module.
\item Let $z=a+bx+cx^{2}$ ($a,b,c\in\mathbf{Q}$) be an integer of $K$. By
calculating the traces of $z,xz$, and $x^{2}z$ show that
$6A\subset B$.
\end{alist}

For the rest of the proof one may, by making use of \hyperref[exer:R2.2]{(2)}, show that the
integers $6a,6b$, and $6c$ defined in (b) are multiples of $6$ (the argument
is long, but elementary). One may also proceed, in a less elementary way, as
follows.
\begin{alist}
\item Show that the ideals $xA$ and $xB$ are prime (for $xA$ one may
consider the decomposition of the ideal $2A$ in $A$). Show that
$xB=B\cap xA$, and that the fields $A/xA$ and $B/xB$ are the same. Conclude
that $A=B+xA$ and, hence, that $A=B+2A$.
\item Show that $3=(x-1)(x+1)^{3}$ and that $x-1$ is invertible in $B$.
Proceed as in (c), with $x+1$ in place of $x$ and $3$ in place of $2$, to
show that $A=B+(x+1)A$, and consequently $A=B+3A$.
\item Deduce from (c) and (d) that $A=B+6A$, and, finally, deduce that
$A=B$.
\end{alist}

\exitem{4} Calculate the different and the discriminant of $K$.

\exitem{5} (a) What is the structure of the (multiplicative) group of units
of $A$? Show that every positive unit (respectively, negative unit) is of
norm $1$ (respectively, $-1$). Show that $x^{2}+x+1$ is a unit of $A$.
\begin{alist}
\item Recall that, if $u$ is a unit $>1$ of $A$ and if $D$ is the
discriminant of $K$, then $|D|<4u^{3}+24$ (cf.\ \hyperref[exer:IV.3]{Exercise 3} of \hyperref[chap:IV]{Chapter IV}).
Conclude that $x^{2}+x+1$ is a generator of the group of positive units of
$K$ (the approximate values of $x$ and $x^{2}$ are $1.26$ and $1.58$).
\item Calculate $3$ or $4$ solutions in integers $a,b,c$ of the equation
$a^{3}+2b^{3}+4c^{3}-6abc=1$. Explain how one may find all its rational
integer solutions, also all its positive integer solutions.
\end{alist}

\exchap{R3}

\subsechead{}{III}

\exitem{1} Show that the polynomial $X^{3}-3X+1$ has no root in
$\mathbf{Q}$ and is irreducible over $\mathbf{Q}$. Show that its three roots
in $\mathbf{C}$ are real; let $x$ be one of them. Write $K$ for the cubic
field $\mathbf{Q}[x]$ and $A$ for the ring of integers of $K$.

\exitem{2} Calculate the discriminant $D(1,x,x^{2})$. Deduce that, if
$a+bx+cx^{2}$ ($a,b,c\in\mathbf{Q}$) is an integer of $K$, then
$a,b,c\in S^{-1}\mathbf{Z}$ where $S$ is the set of powers of $3$,
$\{3^{n}\}_{n\geq0}$ (consider $S^{-1}\mathbf{Z}\subset S^{-1}A$).

\exitem{3} Put $A'=\mathbf{Z}[x]$. Show that $x$ and $x+2$ are units of
$A'$ and of $A$, and that $(x+1)^{3}=3x(x+2)$. Conclude that $(x+1)A$ is a
prime ideal of $A$, that $(x+1)A\cap A'=(x+1)A'$, and that
$A/(x+1)A=A'/(x+1)A'$. Show that $A=A'+(x+1)A$ and $A=A'+3A$. Conclude from
this and from \hyperref[exer:R3.2]{(2)} that $A=A'=\mathbf{Z}[x]$.

\exitem{4} Show that $2A$ is a prime ideal of $A$ (note that
$X^{3}-3X+1$ is irreducible mod $2$). Deduce from this that $A$ is a
principal ideal ring (recall that every ideal class of $A$ contains an
integral ideal of norm
$\leq(4/\pi)^{r_{2}}(n!/n^{n})|D|^{1/2}$, with the usual notations).

\exitem{5} Show that $x^{2}-2$ is a root of the polynomial $X^{3}-3X+1$.
What are the fields conjugate to $K$ over $\mathbf{Q}$, and what are the
$\mathbf{Q}$-isomorphisms of $K$ into $\mathbf{C}$?

\exitem{6} Recall that $\cos3u=4\cos^{3}u-3\cos u$. Putting $x=2\cos u$,
calculate the possible values for the angle $u$. Conclude that $K$ is of the
form $L\cap\mathbf{R}$, where $L$ is a cyclotomic field. Determine which
cyclotomic field.

\enlargethispage{2\baselineskip}
\medskip
\noindent N.B. Questions 5 and 6 are independent of questions 2, 3, and 4.
